\documentclass[11pt, a4paper]{article}
\usepackage[utf8]{inputenc}
\usepackage[ngerman]{babel}
\usepackage[T1]{fontenc}
\usepackage{geometry}
\geometry{a4paper, left=25mm, right=25mm, top=25mm, bottom=25mm}
\usepackage{helvet}
\renewcommand{\familydefault}{\sfdefault}
\usepackage{enumitem}
\usepackage{titlesec}
\usepackage{hyperref}
\usepackage{xcolor}
\usepackage{amssymb}
\usepackage{array}
\usepackage{booktabs}
\usepackage{listings}
\usepackage{mdframed}

\titleformat{\section}{\large\bfseries}{}{0em}{}[\vspace{-0.3em}]
\titleformat{\subsection}{\normalsize\bfseries}{}{0em}{}
\titleformat{\subsubsection}{\normalsize\itshape}{}{0em}{}

\setlength{\parindent}{0pt}
\setlength{\parskip}{0.5em}

\definecolor{linkcolor}{RGB}{0,70,140}
\hypersetup{colorlinks=true, linkcolor=linkcolor, urlcolor=linkcolor}

\lstset{
  basicstyle=\ttfamily\small,
  backgroundcolor=\color{gray!10},
  frame=single,
  breaklines=true
}

\begin{document}

% ============================================================
% HEADER
% ============================================================
\begin{center}
    {\large\bfseries Master-Dokumentation: 3DGS Scan-to-BIM Pipeline}\\[0.5em]
    {\normalsize\bfseries KI-gest\"utzte Vorverarbeitung, Machbarkeitsanalyse und Docker-Infrastruktur}\\[0.8em]
    \small
    \begin{tabular}{ll}
        \textbf{Studiengang:} & Geovisualisierung (B.Eng.) \\
        \textbf{Datum:}       & \today                     \\
    \end{tabular}
\end{center}

\vspace{-0.5em}
\noindent\rule{\textwidth}{0.4pt}
\tableofcontents
\vspace{1em}
\noindent\rule{\textwidth}{0.4pt}

% ============================================================
\section{1. Die aktuelle Pipeline-Idee (Der modifizierte 3-Schritte-Plan)}

Das Ziel dieser Pipeline ist es, schwer zu erfassende, oft extrem d\"unne Objekte (wie Stromleitungen oder gest\"uckelte Schl\"auche) pr\"azise zu isolieren und als sauberes 3D-Gittermodell (Mesh) zu rekonstruieren. Der Workflow basiert auf einer klugen Entkopplung von 2D-Segmentierung und 3D-Training.

\textbf{Die modifizierte Architektur:}
\begin{enumerate}[noitemsep]
    \item \textbf{Das COLMAP-Fundament:} Unbearbeitete Drohnenbilder werden in COLMAP verarbeitet, um das Kamera-Tracking durch den detailreichen Hintergrund abzusichern.
    \item \textbf{Intelligente 2D-Vorarbeit (Schritt 1):} Grounded-SAM (textbasiert) und DEVA (Videotracking) weisen den Kabel-Pixeln in den Bildern ID-Masken zu.
    \item \textbf{Logisches B\"undeln (Schritt 2):} Ein Python-Skript f\"arbt den Hintergrund aller Bilder rabenschwarz.
    \item \textbf{Der Bait \& Switch (Schritt 3):} Die Originalbilder im COLMAP-Ordner werden durch die schwarzen Bilder ausgetauscht.
    \item \textbf{Fokussiertes 3D-Meshing (Schritt 4):} SuGaR trainiert 100\% seiner 3D-Gaussians ausschlie\ss lich auf den Leitungen und extrahiert daraus das finale 3D-Mesh.
\end{enumerate}

% ============================================================
\section{2. Machbarkeitsanalyse (Feasibility Report)}

\subsection{Phase 1: Infrastruktur (WSL2, Docker, CUDA)}
Der native Windows-Betrieb (MSVC Compiler) f\"uhrt bei diesen KI-Repos garantiert zu unl\"osbaren Kompilierungsfehlern. Windows Subsystem for Linux (WSL2) mit Ubuntu 22.04 LTS ist zwingend erforderlich. 
Das Problem der "Dependency Hell" (alte CUDA 11.3 Version f\"ur Vorverarbeitung vs. CUDA 11.8 f\"ur SuGaR) l\"asst sich \textbf{nur durch Docker} l\"osen. Durch die Nutzung des NVIDIA Container Toolkits l\"aufen die Teilschritte in streng isolierten Containern. Die Masken m\"ussen zudem nicht in einem KI-Modell gespeichert werden; sie werden effizient als PNGs in geteilten Ordnern abgelegt.

\subsection{Phase 2: Die 2D-Vorverarbeitung (SAM + DEVA)}
Die "Out-of-the-box" Erkennung extrem d\"unner Leitungen scheitert oft. Grounding DINO spannt riesige Bounding Boxes um diagonale Leitungen, woraufhin SAM den Fokus verliert. DEVA verliert das Tracking bei Motion Blur. 
\textbf{Workarounds:} Die Bilder m\"ussen mittels SAHI (Tiling) in Kacheln zerschnitten werden. Zur Vermeidung von Motion Blur m\"ussen die Drohnenaufnahmen mit sehr kurzen Belichtungszeiten und hohen Framerates (60+ fps) gemacht werden.

\subsection{Phase 3: Der fatale COLMAP-Flaschenhals}
Ein massiver Architekturfehler im Ursprungsplan: COLMAP st\"urzt ab (Aperture Problem), wenn es freigestellte Bilder (Leitung auf schwarzem Hintergrund) verarbeiten soll, da SIFT keine Hintergrund-Features findet.
\textbf{Die L\"osung:} Der "Image Swap". COLMAP berechnet die Kamerapositionen auf den \textit{Originalbildern}. Erst danach werden die Bilddateien auf der Festplatte gegen die schwarzen Masken ausgetauscht, bevor SuGaR gestartet wird.

\subsection{Phase 4: Das 3D-Meshing (SuGaR vs. Poisson)}
Die in SuGaR verwendete Poisson Surface Reconstruction baut bevorzugt geschlossene Volumen und hasst offene, hauchd\"unne R\"ohren (Gefahr des Oversmoothing). Da jedoch 100\% der Gaussians (durch den schwarzen Hintergrund) auf die Leitung gezwungen werden, entsteht eine extreme Punktdichte, die Poisson stabilisiert. 
\textbf{Plan B:} Falls das Mesh dennoch rei\ss t, k\"onnen die Zentren der Gaussians direkt als dichte Punktwolke (.ply) an Software wie DGtal \"ubergeben werden, um perfekte CAD-Zylinder zu fitten.

% ============================================================
\section{3. Praxis-Handbuch: Docker \& System-Setup}

Sollten Sie das Projekt auf Ihrem Windows-Laptop via WSL2 machen, oder sich einen echten Linux-Rechner (z.B. im Uni-Labor) besorgen? Ein nativer Linux-PC hat direkten GPU-Zugriff und verursacht weniger Kopfschmerzen. M\"ussen Sie unter Windows arbeiten, ist WSL2 der einzige Weg. Wichtig: Speichern Sie Ihre Daten niemals auf den C:/ Laufwerken, sondern zwingend innerhalb des virtuellen Linux-Dateisystems (\texttt{\textbackslash\textbackslash wsl\$\textbackslash Ubuntu\textbackslash home\textbackslash...}).

\subsection{Administrator-Rechte}
Admin-Rechte unter Windows werden ben\"otigt f\"ur:
\begin{itemize}[noitemsep]
    \item Die WSL2 Aktivierung (\texttt{wsl -{}-install}).
    \item Die Installation von Docker Desktop f\"ur Windows (mit WSL2-Backend).
\end{itemize}
NVIDIA Treiber m\"ussen nicht im Linux installiert werden, das \"ubernimmt Windows automatisch.

\subsection{Abh\"angigkeiten \& Issue \#20 (DEVA)}
Ein ber\"uhmter C++ Compiler-Fehler (Issue \#20) im DEVA-Repo unter WSL2 tritt auf, weil sich die Windows- und Linux-Umgebungsvariablen verheddern. Durch die Nutzung von Docker wird dieser Fehler obsolet, da der Container (z.B. \texttt{nvidia/cuda:11.7.1-devel-ubuntu22.04}) seinen eigenen, sauberen \texttt{nvcc} Compiler mitbringt.

\subsection{Die Docker-Compose Architektur (Bauplan)}
Die Trennung der Pipeline erfolgt am besten \"uber eine \texttt{docker-compose.yml}:

\begin{lstlisting}
version: '3.8'
services:
  vorverarbeitung:
    build:
      context: ./Repos_Pipeline/Tracking-Anything-with-DEVA
    volumes:
      - ./mein_drohnen_datensatz:/daten_ordner
    deploy:
      resources:
        reservations:
          devices:
            - driver: nvidia
              count: all
              capabilities: [gpu]
    command: /bin/bash

  meshing:
    build:
      context: ./Repos_Pipeline/SuGaR
    volumes:
      - ./mein_drohnen_datensatz:/daten_ordner
    deploy:
      resources:
        reservations:
          devices:
            - driver: nvidia
              count: all
              capabilities: [gpu]
    command: /bin/bash
\end{lstlisting}

% ============================================================
\section{4. Repository-Analysen}

\begin{itemize}
    \item \textbf{lkeab/gaussian-grouping:} Genutzt als 2D-Vorverarbeitungs-Werkzeug (Schritt 1) zur Erstellung konsistenter Masken via SAM+DEVA, ohne das 3D-Training zu nutzen.
    \item \textbf{hkchengrex/Tracking-Anything-with-DEVA:} Das Herzst\"uck des Frame-Trackings. Trennt Bildsegmentierung vom zeitlichen Tracking.
    \item \textbf{hkchengrex/Grounded-Segment-Anything:} Kombination aus Zero-Shot Objekterkennung (DINO) und präziser Segmentierung (SAM) per Text-Prompt (z.B. "power line").
    \item \textbf{Anttwo/SuGaR:} Der "Finisher" (Schritt 4). Extrahiert ein hochpr\"azises 3D-Mesh aus den flach an der Oberfl\"ache anliegenden Gaussians.
    \item \textbf{florinshen/FlashSplat:} Optionaler Plan-B Benchmark f\"ur extrem schnelle 2D-zu-3D Projektion.
\end{itemize}

\end{document}
